\bookchapter{Capítulo 28: Gastar bien para comprar el móvil ideal — Guía definitiva de compra}

\chapterintro{Comprar un teléfono: mucha gente mira reseñas durante una semana y compara precios durante un mes, pero al final gasta cinco o seis mil yuanes en un teléfono insignia que en realidad no necesita. Este capítulo no lista tablas de parámetros; eso lo puedes ver en la web oficial. Lo que quiero darte es un “método de pensamiento para elegir teléfono”, para que en el futuro, sin importar qué teléfono nuevo salga al mercado, puedas juzgar por ti mismo si te conviene.}

\booksection{Método de tres pasos: necesidades presupuesto modelo}

Mucha gente pregunta en internet “¿qué móvil comprar por 3000 yuanes?”, pero esta pregunta no se puede responder sin saber para qué vas a usar el teléfono. El orden correcto de pensamiento al comprar debería ser:

\begin{tcbinfo}{Tres pasos para elegir teléfono}
\begin{enumerate}
  \item \textbf{Primero piensa qué es lo que más te importa}: ¿Fotos? ¿Juegos? ¿Autonomía? ¿Pantalla? ¿Fluidez del sistema? ¿O solo necesitas un teléfono que pueda hacer llamadas y usar WeChat? Ordena tus necesidades.
  \item \textbf{Luego fija un presupuesto máximo}: No "unos cuántos miles", sino "cuánto puedo gastar como máximo". Fija un límite firme para no ir viendo modelos cada vez más caros.
  \item \textbf{Solo al final elige el modelo}: Con tu lista de necesidades ordenadas y tu presupuesto, busca modelos que coincidan.
\end{enumerate}

\end{tcbinfo}

\booksection{I. Cómo leer los parámetros clave - no necesitas entenderlos todos, con estos cinco basta}

Las fichas técnicas tienen decenas de datos, pero los usuarios comunes solo necesitan prestar atención a cinco dimensiones:

\begin{center}
\begin{tabularx}{\linewidth}{|X|X|X|}
\hline
\textbf{Dimension} & \textbf{Cómo mirarlo} & \textbf{Resumen en una frase} \tabularnewline
\hline
Procesador (SoC) & Determina cuántos años el teléfono funcionará fluido. Estándar 2026: Snapdragon 8 Elite Gen5/Dimensity 9500 o superior = gama alta; Snapdragon 8s Gen4/Dimensity 8400 = gama media-alta; Snapdragon 7+ Gen4/Dimensity 7350 = suficiente para tres años. Por debajo de Snapdragon 6 series, considérelo con cuidado. & Compra lo nuevo no lo viejo; el SoC determina cuántos años tu teléfono funcionará sin lag \tabularnewline
Pantalla & OLED tiene mejor color y consume menos que LCD; 120Hz es mucho más fluido al deslizar que 60Hz; resolución mínima 1080P. El PWM de alta frecuencia (1920Hz+) es mejor para los ojos. & OLED + 120Hz + PWM de alta frecuencia = línea de aprobación de buena pantalla \tabularnewline
Android & \textbf{No mires los megapíxeles, mira el tamaño del sensor}. Un sensor más grande es mejor; el tamaño del sensor es mucho más importante que los megapíxeles. El gran angular es práctico; el teleobjetivo (3x o más zoom óptico) para fotografiar cosas lejanas; la lente macro casi siempre es de relleno. & Sensor principal grande + teleobjetivo práctico > múltiples lentes de relleno \tabularnewline
Batería y carga & 4500mAh es el mínimo aceptable; 5000mAh+ para un día sin preocupaciones. La carga rápida se mide en vatios: más de 60W es rápida; más de 100W es rapidísima. La carga inalámbrica es cómoda pero lenta. & 5000mAh + carga rápida 60W = combinación sin ansiedad \tabularnewline
Compromiso de actualización del sistema & Determina cuánto tiempo el teléfono recibirá actualizaciones de seguridad. Apple suele dar 5-6 años de actualizaciones; Samsung insignia da 4 años de sistema + 5 años de seguridad; Google Pixel da 7 años; los gama media nacionales tienen compromisos más cortos (2-3 años). & Si quieres usarlo mucho tiempo, elige una marca con largo compromiso de actualizaciones \tabularnewline
\hline
\end{tabularx}
\end{center}

\booksection{II. Recomendaciones por rango de presupuesto (basado en el mercado chino a mediados de 2026)}

Las siguientes recomendaciones se basan en el mercado de China continental en 2026; los precios son oficiales de salida. Los precios concretos pueden variar; usa el precio real al momento de compra.

\booksubsection{Menos de 1000 yuanes: llamadas, WeChat, suficiente}

\textbf{¿Para quién?}: Estudiantes, teléfono de respaldo, usuarios con pocos requisitos, principalmente para llamadas y WeChat.

\textbf{¿Qué puedes comprar?}: Redmi Note series, Honor X series, realme serie Digital de gama baja. Procesador Snapdragon 6 series o Dimensity 6000; pantalla principalmente LCD; fotos decentes pero enfoque lento; fluidez general.

\textbf{Recomendacion de compra}: En este rango, mira dos cosas: si la batería es suficientemente grande (elige 5000mAh+) y si hay suficiente memoria (mínimo 6GB+128GB, no compres 4GB porque irá lento). Redmi tiene la mejor ventaja en este rango.

\booksubsection{1000-2000 yuanes: la mejor relación calidad-precio, zona dulce para usar 3-4 años sin lag}

\textbf{¿Para quién?}: Estudiantes, presupuesto limitado pero quieren durabilidad, sin exigencias extremas en fotos o juegos.

\textbf{Modelos representativos}: Redmi Turbo series, iQOO Z series, Honor X series gama alta, OnePlus Ace series entrada.

\textbf{Recomendacion de compra}: Elige modelos con Snapdragon 7+ Gen3 o Dimensity 7300+; con 8GB+256GB. Esta configuración es suficiente para uso diario + ocasional Honor of Kings; puede durar 3-4 años sin lag.

\booksubsection{2000-3500 yuanes: gama media, rendimiento y fotografía con mejora notable}

\textbf{¿Para quién?}: Usuarios que quieren buenas fotos, juegos sin lag y durabilidad. Este es el rango de precio donde más gente compra.

\textbf{Modelos representativos}: Xiaomi serie Digital (como Xiaomi 17 standard tras bajada de precio), vivo S series gama alta, OPPO Reno series, Honor serie Digital, Redmi K series.

\textbf{Recomendacion de compra}: La característica principal de este rango es ser “telércubo” sin puntos débiles evidentes. Frente a los de menos de 2000 yuanes, la mejora principal está en: cámara notablemente mejor (con estabilización óptica), pantalla OLED de alta calidad, carga rápida (80W+ habitual), mejor construcción y acabados. Elige 12GB+256GB o 12GB+512GB, usable 4-5 años sin problemas.

\booksubsection{3500-5000 yuanes: cuasi-insignia, suficiente para todos excepto geeks}

\textbf{¿Para quién?}: Usuarios que buscan buena fotografía, quieren experimentar calidad premium.

\textbf{Modelos representativos}: Xiaomi serie Digital standard, vivo X series, OPPO Find X series standard, iPhone antiguo insignia (como iPhone 16).

\textbf{Recomendacion de compra}: En este rango ya puedes comprar el teléfono insignia del año pasado. Mi consejo es:\textbf{compra el insignia del año pasado, no el gama media de este año}. Los materiales, construcción y hardware de cámara del insignia del año pasado suelen ser un nivel superior al gama media de este año. La premisa es que aceptes que no es “el último modelo”.

\booksubsection{Más de 5000 yuanes: zona insignia, dinero no es problema o búsqueda de experiencia extrema}

\textbf{¿Para quién?}: Entusiastas de la fotografía, gamers intensivos, usuarios que buscan pantalla y artesanía premium, creadores de contenido que graban video con el teléfono.

\textbf{Modelos representativos}: iPhone 17 Pro/Pro Max, Samsung S26 Ultra, Huawei Pura 80 Ultra, Xiaomi 17 Ultra, vivo X300 Pro.

\textbf{Recomendacion de compra}: En este rango, las diferencias entre marcas son cada vez menores; principalmente depende de qué sistema te guste. Si te gustan las fotos, elige vivo o Huawei; si te gusta grabar video, elige iPhone; si te gustan la pantalla y el S Pen, elige Samsung; si buscas equilibrio, elige Xiaomi. Sinceramente, a este precio,\textbf{puedes ser caprichoso y elegir el que te guste visualmente}.

\booksection{III. Elección por escenario - ¿cuál es tu necesidad?}

Estos son algunos de los “perfiles de uso” más comunes; encuéntrate en ellos:

\begin{center}
\begin{tabularx}{\linewidth}{|X|X|X|}
\hline
\textbf{Escenarios de uso} & \textbf{Qué mirar} & \textbf{Dirección recomendada} \tabularnewline
\hline
Entusiasta de la fotografía & Sensor principal grande (1 pulgada ideal), estabilización óptica, teleobjetivo (3x+ zoom óptico), chip de imagen & vivo X series, Huawei Pura/Mate series, Xiaomi Ultra series \tabularnewline
Gamer & Procesador tope, alta tasa de muestreo táctil, sistema de refrigeración, batería grande + carga rápida & iQOO serie Digital, ROG Phone, Red Magic, Redmi K series \tabularnewline
Autonomía prioritaria & Batería de 5500mAh+, procesador eficiente, pantalla LTPO & Honor Magic series, OnePlus, Redmi Note versión gran batería \tabularnewline
Para mayores & Modo simple, fuente grande, volumen alto, batería grande, muchos puntos de servicio postventa & Huawei Enjoy, Honor Play, OPPO A series, Redmi Note \tabularnewline
Para estudiantes & Alta relación calidad-precio, duradero, batería grande, con control parental & Redmi Note/Turbo, realme GT, iQOO Neo \tabularnewline
Usuario del ecosistema Apple & Ya tienes Mac/iPad/AirPods, comprar iPhone es la única opción & iPhone 17 standard (relación calidad-precio) o Pro (experiencia completa) \tabularnewline
Viajes internacionales & Soporta eSIM, cobertura completa de bandas globales, servicios de Google & iPhone, Samsung S series, Pixel, OnePlus \tabularnewline
\hline
\end{tabularx}
\end{center}

\booksection{IV. Cinco trampas comunes - revísalas antes de comprar}

\begin{enumerate}
  \item \textbf{Trampa 1: cuatro o cinco cámaras no significa buenas fotos}; los fabricantes acumulan lentes en las especificaciones, pero a menudo dos de ellas son de 2 megapíxeles y son de "relleno" (macro + profundidad), completamente inútiles en la práctica. Para evaluar la capacidad fotográfica de un teléfono,\textbf{solo mira la cámara principal + gran angular + teleobjetivo}. El resto puedes ignorarlo.
  \item \textbf{Trampa 2: no te dejes impresionar por 100 megapíxeles}; muchos píxeles no significa buena calidad. El tamaño del sensor, la apertura y el procesamiento algorítmico son mucho más importantes. Muchos teléfonos insignia usan sensores grandes de 50MP y toman fotos mucho mejores que algunos teléfonos de 100MP de mil yuanes.
  \item \textbf{Trampa 3: no compres teléfonos personalizados de operador}; "recarga y te llevas un teléfono" parece muy rentable, pero los teléfonos personalizados de operador tienen actualizaciones lentas, pueden tener apps preinstaladas que no se pueden eliminar y suelen estar bloqueados por el operador (solo funcionan con la SIM de ese operador). Parece que ahorras dinero, pero en realidad te "bloquea".
  \item \textbf{Trampa 4: "las pantallas curvas son muy premium" - es marketing, no verdad}; las pantallas curvas solían ser símbolo de gama alta, pero cada vez más gente descubre sus problemas: difícil poner protector, toques accidentales en los bordes, reparación cara si se rompe. La tendencia de 2026 es volver a "micro-curvo" o incluso "pantalla plana". No elijas pantalla curva solo porque "lo curvo es más premium".
  \item \textbf{Trampa 5: no todos los "AI phones" son realmente útiles}; en 2026 todos los fabricantes gritan "AI phone", pero muchos solo han cambiado el nombre al asistente de voz que ya tenían. Las funciones de IA realmente útiles son: traducción en tiempo real, eliminación con IA (quitar personas de las fotos), edición de fotos con IA, inferencia de modelos grandes en el dispositivo, resúmenes inteligentes localizados. Las demás funciones vistosas son básicamente trucos de marketing.
\end{enumerate}

\booksection{V. ¿Cuándo es más barato comprar?}

\begin{itemize}
  \item \textbf{618 (finales de mayo a mediados de junio)}: el periodo con mayores descuentos del año. Muchos fabricantes reducen el precio de los insignia entre 800-1500 yuanes para impulsar ventas. Es la mejor ventana para comprar en la primera mitad del año.
  \item \textbf{Doble 11 (finales de octubre a mediados de noviembre)}: la gran promo de la segunda mitad del año, con descuentos similares al 618. En este periodo, muchas marcas lanzan nuevos insignia y los de la generación anterior bajan más de precio.
  \item \textbf{Un mes después del lanzamiento}: al lanzamiento el precio es de lista; un mes después, el precio de canal (tiendas de terceros) suele ser 200-500 yuanes más barato.
  \item \textbf{Dos semanas antes del Año Nuevo Chino hasta el Festival de los Faroles}: las tiendas físicas liquidan stock en la "Feria de Año Nuevo"; se puede regatear con el dependiente.
  \item \textbf{Momentos no recomendados}: comprar el primer día tras la presentación del producto; a menos que realmente quieras ser de los primeros, el precio de lanzamiento es el más caro.
\end{itemize}

\booksection{VI. ¿Dónde comprar? - ¿qué canal es fiable?}

\begin{itemize}
  \item \textbf{Tienda oficial de la marca / flagship store de Tmall / JD.com self-operated}: lo más seguro; sin ventaja de precio pero sin preocuparse por refabricados ni falsificaciones; mejor servicio postventa.
  \item \textbf{Pinduoduo "subsidio de 10.000 millones"}: el precio más bajo, pero elige tiendas marcadas con "subsidio de 10.000 millones". Las no marcadas no las toques. Recuerda grabar un video de desempaquetado; si hay problemas, puedes devolverlo.
  \item \textbf{Tienda física}: puedes tocar el teléfono y sentirlo. Atención: los dependientes recomiendan los modelos con mayor margen (no los más adecuados para ti); te recomiendo que antes de entrar decidas 2-3 modelos objetivo y al entrar digas "solo quiero ver estos".
  \item \textbf{Plataformas de segunda mano (Xianyu / Zhuanzhuan)}: más barato pero con riesgos. Si compras de segunda mano, recuerda tres cosas: pide factura de compra o captura del pedido, pide el número IMEI para verificar,\textbf{solo opera a través de la plataforma, nunca transfieras dinero privadamente}.
\end{itemize}

\begin{tcbpractice}{Práctica: usa este método para elegir un teléfono}
\begin{enumerate}
  \item Coge un papel y escribe tus dos necesidades principales (por ejemplo: autonomía primero, fotos segundo).
  \item Fija un presupuesto máximo (por ejemplo, 3500 yuanes).
  \item Abre JD.com o Tmall, filtra por precio y marca 3-5 modelos candidatos.
  \item Busca en Bilibili "modelo XX - un mes de uso" (no reseñas del primer día, sino experiencias reales tras al menos un mes); mira dos o tres videos.
  \item Confirma que el canal es fiable y haz el pedido; graba un video al desempaquetar.
\end{enumerate}

\end{tcbpractice}

\begin{tcbsummary}{Resumen del capítulo}
\begin{itemize}
  \item Tres pasos para elegir teléfono: primero necesidades luego presupuesto finalmente modelo
  \item Parámetros clave solo cinco: procesador, pantalla, cámara, batería y carga, compromiso de actualización del sistema
  \item 1500-3500 yuanes es la zona dulce con mejor relación calidad-precio
  \item Comprar el insignia del año pasado es mejor que el gama media de este año
  \item Cinco trampas: estafa de múltiples cámaras, trampa de alto megapíxel, personalización de operador, mito de pantalla curva, truco de IA
  \item 618 y Doble 11 son las dos mejores ventanas del año para comprar teléfono
\end{itemize}

\end{tcbsummary}
